
\subsubsection{6.1 Interpretations of β₄ ≈ 0}

Three interpretations are consistent with the null result on β₄; they are not mutually exclusive.

\textbf{Split invalidity.} The most parsimonious account, and the one most directly supported by dataset documentation, is that the HH-RLHF helpful-base and helpful-online splits do not encode annotator exposure history. The split boundary is defined by data collection period and response generation provenance — specifically, which model produced the responses annotators evaluated — not by which individual annotators participated or what their prior AI interaction histories were. If both splits contain comparable proportions of AI-familiar and AI-naive annotators (as is plausible given the gradual diffusion of AI product adoption across the general population), the split variable would be a noisy and possibly uncorrelated proxy for the exposure construct regardless of whether the true adaptation effect exists. Under this interpretation, the null β₄ reflects measurement failure rather than an empirical test of the adaptation hypothesis.

\textbf{Societal saturation.} An alternative interpretation is that, by the time the HH-RLHF annotations were collected, AIFS-style norms had permeated sufficiently broadly — through LLM-generated content appearing in everyday digital contexts — that no meaningful contrast between AI-familiar and AI-naive annotator populations existed. Under this account, the bidirectional adaptation effect may have already become uniform across the population prior to dataset collection, making between-condition detection impossible regardless of how annotator conditions were operationalized.

\textbf{True null at AIFS level.} A third interpretation is that the adaptation hypothesis is empirically incorrect at the level of AIFS surface features: annotator AI familiarity does not modulate preference for structured lists, safety prefaces, chain-of-thought markers, or hedging language, because these features are preferred by all annotators for reasons independent of AI exposure — for example, because structured formatting aids readability regardless of its connotations. The positive β₁ is compatible with this interpretation: AIFS features may be universally useful formatting signals rather than culturally marked indicators of AI-generated content.

The available data do not permit discrimination among these three interpretations, because all three predict β₄ ≈ 0 under the observational design used. Distinguishing them requires either longitudinal annotator-level data (to test interpretation 1) or pre-AI-era preference data as a reference condition (to test interpretation 2).

\subsubsection{6.2 AIFS Construct Validity}

The positive and significant β₁ = +0.0246 (p < 0.001) is the principal affirmative finding of this study. It demonstrates that the AIFS composite — four regex-based pattern counts — captures preference-relevant variation that survives controls for response length and cluster fixed effects. The effect size is modest (OR ≈ 1.025 per unit of ΔAIFS), consistent with AIFS being one signal among multiple contributors to annotator preference rather than a dominant factor.

The construct validity of AIFS is conditional on the representativeness of the HH-RLHF corpus. The AIFS patterns were defined by inspection of RLHF-trained model outputs and have not been validated against pre-LLM writing baselines to confirm that they are differentially present in AI-generated versus human-written text. Future work should assess whether AIFS scores discriminate between AI-generated and human-written responses from the same period, which would strengthen the interpretive basis for the construct.

\subsubsection{6.3 Limitations}

\textbf{Single dataset.} All analyses use HH-RLHF, a single dataset from one organization with a specific annotation protocol. Whether AIFS features predict preference in other corpora — OpenAssistant, UltraFeedback, LMSYS-Chat-1M — has not been tested.

\textbf{No annotator identity records.} The HH-RLHF public release does not include annotator identifiers, demographic information, or AI usage histories. Annotator pool overlap across splits cannot be assessed. Verification of the proxy validity assumption is therefore not possible within this dataset.

\textbf{AIFS not validated against pre-LLM baselines.} The four feature classes were selected based on known properties of RLHF-trained outputs, but empirical discrimination from human-written text has not been established.

\textbf{Societal saturation untestable.} Whether "naive" annotators in the base split had meaningful AI exposure at the time of data collection cannot be determined from available information. This limits the interpretive options available for the null result.

\textbf{Cluster independence assumption.} The precision estimates in Section 5.3 assume independence across the 27,034 clusters. Cross-cluster dependencies arising from shared annotator pools cannot be ruled out in the absence of annotator metadata.

\textbf{Blocked sub-hypotheses.} Four mechanism sub-hypotheses (H-M1 through H-M4) were planned but were not executed because the existence-level hypothesis (H-E1) failed its MUST_WORK gate. Results for these sub-hypotheses are therefore not available.

\subsubsection{6.4 Data Infrastructure Requirements}

The present study illustrates a structural requirement for empirical bidirectional alignment research: preference datasets must include annotator-level longitudinal exposure metadata. The split-based design used here — treating data collection provenance labels as annotator exposure indicators — is not supported by the available evidence and does not yield interpretable results for the adaptation hypothesis.

A dataset design that would support the planned analysis would require: annotator identifiers enabling longitudinal tracking, verified or self-reported AI product usage histories, and annotation collection spanning a meaningful pre-/post-exposure interval. Datasets with relevant partial structure include UltraFeedback (annotator identifiers), LMSYS-Chat-1M (session-level metadata), and AlpacaFarm (explicit annotator demographic information). Within-annotator longitudinal designs — collecting preference annotations from the same individuals before and after AI tool adoption — would provide the causal leverage that cross-sectional designs cannot.

The narrow CI on the null result (width 0.025, excluding OR ≥ 1.0108) constrains the hypothesis space by ruling out practically meaningful effects detectable through the split proxy. This does not constrain what effects might be detectable through a properly designed annotator-linked study.

